% Insertion-ready after q444:complex-stress-obstruction.
% This module proves a new bounded finite-fibre elimination theorem.
% It re-proves, but does not claim novelty for, the signed norm identity
% already present in coordinate_constraints/orthogonality_normal_form.tex.
\subsection{An exact bounded finite coupling fibre for every unrestricted spatial pair}
\label{q444:finite-coupling-section}

Retain the literal matrices $H,q,F$, all eight real coordinates of each
$Z_j\in\mathfrak{su}(3)$, and the original positive coupling $g^2$.
The inner product remains $\langle X,Y\rangle=-\operatorname{tr}(XY)$.
The letter $t$ below is a polynomial indeterminate whose evaluated value is
the original $g^2$; no physical constant or amplitude is set to one.
This subsection constructs the complete fibre over a fixed ordered pair
$(Z_1,Z_2)$ of the four-residual zero set
\eqref{q444:complex-full-residuals}. The ordered pair is not constrained
to a real subalgebra, a block class, or a prescribed phase relation.

\begin{proposition}[An explicit upper bound for the original coupling]
\label{q444:finite-coupling-bound}
For an arbitrary fixed ordered pair with $[Z_1,Z_2]\ne0$, define
\[
 z_0=\|Z_1\|^2+\|Z_2\|^2,\qquad
 z=\sum_{j=1}^2\langle Z_j,\mathcal IZ_j\rangle,\qquad
 d=\|[Z_1,Z_2]\|^2,
\]
\[
 c_* =2+z,\qquad k_*=2z_0(z_0+4),\qquad
 t_*=
 \frac{k_*z+\sqrt{k_*^2z^2+8k_*d c_*^2}}{c_*^2}.
\]
All these numbers are real, with $z\ge0$ and
$z_0,d,c_*,k_*,t_*>0$.
Every full solution over this ordered pair satisfies
\begin{equation}
 \|C\|^2\le z_0+4,\qquad
 0<g^2\le t_*,\qquad
 0<\omega^2\le
 \frac{k_*+2d+t_*(c_*+z/2)}{z_0}.
 \label{q444:finite-coupling-bounds}
\end{equation}
The same bounds apply to both signs of the original frequency.
\end{proposition}
\begin{proof}
We include the signed norm calculation so that the amplitude estimate has
no unproved input from a proposed coordinate restriction. In the full
coordinates \eqref{q444:complex-eight-coordinates}, put $h_X=a_X+2b_X$.
Multiplication of the diagonal and off-diagonal entries gives
\[
 \|X\|^2=\tfrac32a_X^2+\tfrac12h_X^2
       +2(|v_{X1}|^2+|v_{X2}|^2+|z_X|^2).
\]
The scalar identities
\eqref{q444:complex-scalar-diagonal-identities} give
$D_a=4D_z$ and $D_1+D_2=2-4D_z$. Consequently
\begin{align*}
 \|C\|^2-z_0
 &=\tfrac32D_a+\tfrac12(h_C^2-h_{Z_1}^2-h_{Z_2}^2)
                    +2(D_1+D_2+D_z)\\
 &=4+\tfrac12(h_C^2-h_{Z_1}^2-h_{Z_2}^2).
\end{align*}
For completeness, the sign of the last correction follows directly from
the full scalar cross identities. Write
$V_X=(v_{X1},\sqrt2 z_X)$,
$S_1=\|V_{Z_1}\|^2+\|V_{Z_2}\|^2$, and
$H_s=h_{Z_1}^2+h_{Z_2}^2$. Those identities give
$h_C V_C=h_{Z_1}V_{Z_1}+h_{Z_2}V_{Z_2}$, while the diagonal identities
give $\|V_C\|^2=S_1+1-\sqrt{2/3}$.
Cauchy--Schwarz therefore proves
\[
 h_C^2(S_1+1-\sqrt{2/3})\le H_s S_1.
\]
Since $1-\sqrt{2/3}>0$, this implies $h_C^2\le H_s$, also when
$S_1=0$ or $H_s=0$. This proves $\|C\|^2\le z_0+4$.

For arbitrary anti-Hermitian matrices $X,Y$, diagonalize $X$ unitarily,
writing its eigenvalues as $i\lambda_a$ with real $\lambda_a$.
In that same basis,
\[
 \|[X,Y]\|^2
 =\sum_{a,b}(\lambda_a-\lambda_b)^2|Y_{ab}|^2
 \le2\Bigl(\sum_a\lambda_a^2\Bigr)
               \Bigl(\sum_{a,b}|Y_{ab}|^2\Bigr)
 =2\|X\|^2\|Y\|^2.
\]
Here $(\lambda_a-\lambda_b)^2\le2(\lambda_a^2+\lambda_b^2)
\le2\sum_a\lambda_a^2$ proves the inequality term by term.
Thus the retained $k=\sum_j\|[C,Z_j]\|^2$ satisfies $k\le k_*$.
The universal identities \eqref{q444:cx:eq:universal} give
$c=c_*$ and
\[
 (g^2)^2c_*^2\le8kd+2g^2kz
                    \le8k_*d+2g^2k_*z.
\]
The quadratic on moving the right side to the left has positive leading
coefficient and roots
$\bigl(k_*z\pm\sqrt{k_*^2z^2+8k_*d c_*^2}\bigr)/c_*^2$.
Its negative root is strictly negative, so its nonpositive region at
$g^2>0$ is exactly $g^2\le t_*$. Finally the universal spatial trace
identity is
\[
 \omega^2z_0=k+2d+g^2(c_*+z/2).
\]
Substitution of the two upper bounds proves the last assertion.
Nonzero bracket implies $z_0,d>0$; positivity of $\mathcal I$ and the
definitions prove all remaining strict and weak sign assertions.
\end{proof}

We now eliminate the temporal stabilizer component exactly before
clearing any frequency denominator. This preserves the original
stabilizer vector $J=iH$, whose squared norm is $3/4$.
Set
\[
 W=\{X\in\mathfrak{su}(3):\langle J,X\rangle=0\},\qquad
 \Pi X=X-\frac{\langle J,X\rangle}{3/4}J,
 \qquad \mathscr B=[Z_1,Z_2].
\]
For the fixed spatial pair define
\[
 \eta=\langle J,\mathcal KJ\rangle,\qquad
 w=\Pi\mathcal KJ,\qquad \beta=\langle J,\mathscr B\rangle,
\]
and operators and a vector on the original metric subspace $W$ by
\[
 \mathcal I_7=\mathcal I|_W,\qquad
 \mathcal S_7=\Pi\mathcal K|_W-\frac{w\otimes w}{\eta},\qquad
 f_7=\Pi\mathscr B-\frac{\beta}{\eta}w.
\]
Here $(w\otimes w)x=w\langle w,x\rangle$; no choice of basis
alters any of these maps.

\begin{lemma}[Full stabilizer elimination and polynomial Gauss numerator]
\label{q444:finite-coupling-schur}
If $\mathscr B\ne0$, then $\eta>0$, $\mathcal S_7\succeq0$, and
$\mathcal I_7\succ0$. Define the polynomial endomorphism and its
determinant and adjugate by
\[
 \mathscr M(t)=t\mathcal I_7+2\mathcal S_7,\qquad
 D_7(t)=\det_W\mathscr M(t),\qquad
 V_7(t)=4\operatorname{adj}_W(\mathscr M(t))f_7,
\]
\begin{equation}
 N_7(t)=V_7(t)+
       \frac{2\beta D_7(t)-\langle w,V_7(t)\rangle}{\eta}J.
 \label{q444:finite-coupling-N}
\end{equation}
Then $D_7(t)>0$ for $t>0$, $D_7$ has degree exactly seven and leading
coefficient $3/128$, $\deg V_7\le6$, and $\deg N_7\le7$.
For every original $g^2>0$,
\begin{equation}
 T=4(g^2\mathcal I+2\mathcal K)^{-1}\mathscr B
   =\frac{N_7(g^2)}{D_7(g^2)},\qquad
 Tq=\frac{V_7(g^2)q}{D_7(g^2)}.
 \label{q444:finite-coupling-T}
\end{equation}
In particular all temporal coordinates, including the entire stabilizer
coefficient, are reconstructed by the displayed formula.
\end{lemma}
\begin{proof}
The positive form of $\mathcal K$ gives
$\eta=\sum_j\|[Z_j,J]\|^2\ge0$. If it were zero, both $Z_j$ would
commute with $H$. The three distinct eigenvalues $0,\sqrt6/4,-\sqrt6/4$
of $H$ would make both matrices diagonal in its eigenbasis; hence
$[Z_1,Z_2]=0$, a contradiction. Thus $\eta>0$.
The kernel calculation in Lemma~\ref{q444:complex-gauss-inverse} gives
$\ker\mathcal I=\mathbb RJ$. Self-adjointness makes $W$ invariant, and
positivity then makes $\mathcal I_7$ positive definite.

For $x\in W$ and $s\in\mathbb R$, use
$\langle J,\mathcal Kx\rangle=\langle w,x\rangle$ to complete the square:
\[
 \langle x+sJ,\mathcal K(x+sJ)\rangle
 =\langle x,\mathcal S_7x\rangle
   +\eta\left(s+\frac{\langle w,x\rangle}{\eta}\right)^2.
\]
Taking $s=-\langle w,x\rangle/\eta$ and using positivity of
$\mathcal K$ proves $\mathcal S_7\succeq0$.
Thus $\mathscr M(t)$ is positive definite for every $t>0$.

Write the full Gauss unknown as $T=x+sJ$ with $x\in W$, with no
restriction on $s$. Pair
$(t\mathcal I+2\mathcal K)T=4\mathscr B$ with $J$ and then apply $\Pi$.
These operations give respectively
\[
 s=\frac{2\beta-\langle w,x\rangle}{\eta},\qquad
 (t\mathcal I_7+2\mathcal S_7)x=4f_7.
\]
Conversely these two equations recover the $J$ component and the $W$
projection of Gauss; the orthogonal decomposition
$\mathfrak{su}(3)=W\oplus\mathbb RJ$ therefore recovers the entire
equation. Inverting the positive operator on $W$ yields
$x=V_7(t)/D_7(t)$, and the formula for $s$ yields
\eqref{q444:finite-coupling-N}. Since $Jq=0$, it also yields the second
formula in \eqref{q444:finite-coupling-T}.

Determinants and adjugates may be calculated in any real basis of $W$;
similarity preserves their resulting scalar and endomorphism. The
inertia computation \eqref{q444:complex-full-inertia-coordinates} gives
the seven eigenvalues $1/6,1/2,1/2,1/2,1/2,3/2,3/2$ on $W$.
Thus $\det_W\mathcal I_7=(1/6)(1/2)^4(3/2)^2=3/128$.
The determinant consequently has degree seven with this leading
coefficient; its positivity for $t>0$ follows from positive definiteness.
The adjugate of a seven-by-seven matrix with affine entries has entries
of degree at most six, and \eqref{q444:finite-coupling-N} proves the
remaining degree bound.
\end{proof}

Put
\[
 \alpha=\frac14+\|Z_1q\|^2+\|Z_2q\|^2>0,\qquad
 Q_7(t)=\|V_7(t)q\|^2.
\]
The real polynomial $Q_7$ is nonnegative at every real $t$ and has degree
at most twelve. Its identically-zero case has the following finite test,
which does not require computing a polynomial determinant.

\begin{proposition}[Three exact vector tests for an empty Gauss domain]
\label{q444:finite-coupling-domain-test}
For every fixed noncommuting spatial pair, $Q_7$ is identically zero if
and only if the following three vector equalities all hold:
\begin{equation}
 f_7q=0,\qquad (\mathcal S_7f_7)q=0,\qquad
                 (\mathcal S_7^2f_7)q=0.
 \label{q444:finite-coupling-three-tests}
\end{equation}
When they hold, no positive original coupling over that spatial pair
belongs to the original four-residual domain.
\end{proposition}
\begin{proof}
Put $\mathcal N=\{X\in W:Xq=0\}$. In the complete coordinates
\eqref{q444:complex-eight-coordinates}, membership of $W$ sets $h_X=0$;
the first column condition $Xq=0$ sets $a_X=v_{X1}=v_{X2}=0$.
The only remaining entry is the arbitrary complex $z_X$. Thus
$\dim_{\mathbb R}\mathcal N=2$, and the full inertia formula gives
$\mathcal I_7|_{\mathcal N}=(3/2)\operatorname{id}_{\mathcal N}$.
Self-adjointness also makes $\mathcal N^\perp\cap W$ invariant.
Consequently $\mathcal I_7^{-1}y\in\mathcal N$ if and only if
$y\in\mathcal N$, and on this subspace the inverse is multiplication
by $2/3$.

Suppose first that $Q_7$ is identically zero. For every $t>0$ the vector
$x(t)=4(t\mathcal I_7+2\mathcal S_7)^{-1}f_7$ then belongs to
$\mathcal N$. Set $B_7=\mathcal I_7^{-1}$.
Factoring the pencil and using the exact finite inverse identity
$(I+A)^{-1}=I-A+A^2-A^3(I+A)^{-1}$ gives
\[
 x(t)=\frac4tB_7f_7-\frac8{t^2}B_7\mathcal S_7B_7f_7
       +\frac{16}{t^3}B_7\mathcal S_7B_7\mathcal S_7B_7f_7
       +O(t^{-4}).
\]
Here the remainder has bounded norm after multiplication by $t^4$:
for $A=2B_7\mathcal S_7/t$ and $t>4\|B_7\mathcal S_7\|$, the inverse
has norm at most two, since
$\|(I+A)y\|\ge(1-\|A\|)\|y\|\ge\|y\|/2$.
The finite-dimensional subspace $\mathcal N$ is closed. Multiplying the
expansion by $t$ and taking the limit shows $B_7f_7\in\mathcal N$, hence
$f_7\in\mathcal N$. Subtracting the first term, multiplying by $t^2$
and taking the limit shows $B_7\mathcal S_7B_7f_7\in\mathcal N$,
hence $\mathcal S_7f_7\in\mathcal N$. Subtracting the next term and
multiplying by $t^3$ similarly gives
$\mathcal S_7^2f_7\in\mathcal N$. In the latter two steps every occurrence
of $B_7$ on the already established vectors in $\mathcal N$ is exactly
multiplication by $2/3$. These are precisely the three displayed tests.

Conversely suppose the three tests hold. If $f_7=0$, then $x(t)=0$ for
all $t>0$ and the conclusion follows. If $f_7\ne0$ and
$\mathcal S_7f_7$ is dependent on $f_7$, their one-dimensional span is
invariant under $\mathcal S_7$. Otherwise those two vectors span all of
$\mathcal N$, and the third test makes that two-dimensional span
invariant as well. In either case let $V$ be the resulting invariant
span. It is contained in $\mathcal N$, so it is invariant under
$\mathcal I_7$ and therefore under the positive pencil
$t\mathcal I_7+2\mathcal S_7$.
Its restriction to $V$ is injective and, in finite dimension, onto $V$;
the full inverse thus carries $f_7\in V$ back into $V$. Hence
$x(t)q=0$ for all $t>0$, and the polynomial $Q_7$ vanishes on the
entire positive half-line. A polynomial with infinitely many roots is
identically zero. Formula \eqref{q444:finite-coupling-T} then proves
the asserted exclusion from the original domain for every $g^2>0$.
\end{proof}

Define four polynomial residuals with the original
matrix and vector codomains:
\begin{align}
 P_H(t)&=\alpha[N_7,[N_7,H]]
     -Q_7\left(\sum_j[Z_j,[Z_j,H]]+F\right),\notag\\
 P_q(t)&=\alpha N_7^2q+Q_7(2I_3-Z_1^2-Z_2^2)q,\notag\\
 P_1(t)&=\alpha\left(-D_7^2Z_1-2D_7[N_7,Z_2]
                       +[N_7,[N_7,Z_1]]\right)
                  +Q_7\left(\mathcal KZ_1+\frac t2\mathcal IZ_1\right),\notag\\
 P_2(t)&=\alpha\left(-D_7^2Z_2+2D_7[N_7,Z_1]
                       +[N_7,[N_7,Z_2]]\right)
                  +Q_7\left(\mathcal KZ_2+\frac t2\mathcal IZ_2\right).
 \label{q444:finite-coupling-polynomial-residuals}
\end{align}
All polynomials on the right are evaluated at the same $t$.
Their codomains are respectively
$\operatorname{Herm}_3(\mathbb C)_0$, $\mathbb C^3$,
$\mathfrak{su}(3)$, and $\mathfrak{su}(3)$ for real $t$.
Define the single real polynomial
\begin{equation}
 p_Z(t)=\operatorname{tr}(P_H(t)^2)+P_q(t)^*P_q(t)
                  -\operatorname{tr}(P_1(t)^2)-\operatorname{tr}(P_2(t)^2).
 \label{q444:finite-coupling-one-polynomial}
\end{equation}

\begin{theorem}[Complete finite coupling fibre and its exact inverse]
\label{q444:finite-coupling-fibre}
For every ordered pair $Z_1,Z_2\in\mathfrak{su}(3)$ the complete fibre of
the four-residual zero set is as follows.
If $[Z_1,Z_2]=0$, it is empty. Otherwise form the explicitly defined
polynomials above. If $Q_7$ is identically zero, the fibre is empty.
If $Q_7$ is not identically zero, then $p_Z$ is not identically zero,
$\deg p_Z\le28$, and the complete set of admissible couplings is exactly
\begin{equation}
 \left\{t\in(0,t_*]:Q_7(t)>0,\ p_Z(t)=0\right\}.
 \label{q444:finite-coupling-root-set}
\end{equation}
This set has at most fourteen elements. Equivalently, the projection
$\mathcal Z\longrightarrow\mathfrak{su}(3)^2$ that forgets $g^2$ has
fibres of cardinality at most fourteen, including empty fibres.

For each root in \eqref{q444:finite-coupling-root-set} and each
$\varepsilon\in\{+1,-1\}$, the complete original finite-system solution is
\begin{equation}
 g^2=t,\qquad
 \omega=\varepsilon\frac{\sqrt\alpha\,D_7(t)}{\sqrt{Q_7(t)}},\qquad
 C=\varepsilon\frac{\sqrt\alpha\,N_7(t)}{\sqrt{Q_7(t)}},\qquad
 (Z_1,Z_2)=(Z_1,Z_2).
 \label{q444:finite-coupling-reconstruction}
\end{equation}
Thus every fibre of the five-equation solution set over the fixed ordered
spatial pair has at most twenty-eight elements, counting both frequency
signs. No root with $Q_7(t)=0$ is included.
\end{theorem}
\begin{proof}
A zero spatial bracket is impossible by
Lemma~\ref{q444:complex-gauss-inverse}. For a nonzero bracket,
\eqref{q444:finite-coupling-T} gives
$\|Tq\|^2=Q_7(t)/D_7(t)^2$ for every $t>0$.
Thus the exact original domain condition $Tq\ne0$ is precisely
$Q_7(t)>0$. If $Q_7$ is identically zero, the domain has no point over
this pair. On the remaining domain the unchanged frequency is
\[
 \tau=\frac{\alpha D_7(t)^2}{Q_7(t)}.
\]
Substitution in all four full residuals gives, entry by entry,
\begin{equation}
 (R_H,R_q,R_1,R_2)
   =\frac1{Q_7(t)}(P_H(t),P_q(t),P_1(t),P_2(t)).
 \label{q444:finite-coupling-exact-residual-map}
\end{equation}
Every denominator here is strictly positive on the asserted domain.
For real $t$, each summand in \eqref{q444:finite-coupling-one-polynomial}
is a squared norm in its stated real inner-product space. The Hermitian
and anti-Hermitian assertions follow directly from the commutators in
\eqref{q444:finite-coupling-polynomial-residuals}; they do not require a
residual to vanish. Therefore $p_Z(t)\ge0$ for every real $t$, and its
vanishing is equivalent to vanishing of every full polynomial residual.
This proves the exact correspondence of roots and full residual zeros,
before the bound on the number of roots is established.

To prove the degree bounds, write $N_7=V_7+s_7J$, with
$\deg V_7\le6$ and $\deg s_7\le7$. Since $[J,H]=0$ and $Jq=0$,
\[
 [N_7,[N_7,H]]=[V_7,[V_7,H]]+s_7[J,[V_7,H]],
 \qquad
 N_7^2q=V_7^2q+s_7JV_7q.
\]
Their degrees are at most thirteen. Hence $\deg P_H,\deg P_q\le13$.
In the two spatial residuals the terms $D_7^2Z_j$,
$D_7[N_7,Z_k]$, and $[N_7,[N_7,Z_j]]$ have degree at most fourteen;
the terms involving $Q_7\mathcal K$ and $tQ_7\mathcal I$ have degree
at most twelve and thirteen. Thus $\deg P_1,\deg P_2\le14$ and
$\deg p_Z\le28$.

Suppose now that $Q_7$ is not identically zero but $p_Z$ is identically
zero. A nonzero univariate polynomial has only finitely many real
roots, so there exists $t>t_*$ with $Q_7(t)\ne0$. Nonnegativity of
$Q_7$ gives $Q_7(t)>0$. The established exact residual correspondence
then makes this $t$ a full zero; the two-sign inverse of
Theorem~\ref{q444:complex-exact-zero-set} reconstructs a solution of all
five original equations. This contradicts
Proposition~\ref{q444:finite-coupling-bound}. Consequently $p_Z$ is a
nonzero polynomial. Every real root of a real polynomial that is
nonnegative on all of $\mathbb R$ has even multiplicity: at a root
$t_0$, factor it as $(t-t_0)^m u(t)$ with $u(t_0)\ne0$; continuity
makes $u$ have a constant nonzero sign near $t_0$, and an odd $m$ would
give opposite signs on the two sides. Thus a nonzero polynomial of
degree at most twenty-eight of this type has at most fourteen distinct
real roots. Imposing the explicit interval and denominator conditions
can only decrease that number.

Finally $D_7(t)>0$ identifies the sign of the displayed $\omega$ with
$\varepsilon$ and gives $C=\omega T$ exactly. The frequency square is
the unchanged $\tau$. Formula
\eqref{q444:finite-coupling-exact-residual-map} proves that all remaining
equations vanish, while the Schur inverse proves the full Gauss equation.
Conversely every original solution determines its fixed pair, its root
$t=g^2$, and its unique sign, and these formulas recover $C,\omega$.
This proves the stated bijection and both cardinality bounds.
\end{proof}

The original spacetime reconstruction is retained literally: with the
already specified $O(t_{\rm phys})$, $K$, and $\Omega$, put
\[
 Y_1(t_{\rm phys})=\cos(\omega t_{\rm phys})Z_1
                         -\sin(\omega t_{\rm phys})Z_2,\qquad
 Y_2(t_{\rm phys})=\sin(\omega t_{\rm phys})Z_1
                         +\cos(\omega t_{\rm phys})Z_2,\qquad Y_3=0,
\]
\[
 \mathcal A_0=OKCK^*O^*-\Omega,\qquad
 \mathcal A_i=OKY_iK^*O^*.
\]
The prescribed scalar fields remain
$\Phi(t_{\rm phys})=O\Phi_0O^*$ and $b(t_{\rm phys})=Ob_0$, with their
original frequency $\sqrt2$. The polynomial indeterminate $t=g^2$ is
distinct from the physical time $t_{\rm phys}$.

This theorem solves the coupling fibre exactly as a bounded finite root
set while keeping every unrestricted spatial coordinate. It does not
establish that any ordered spatial pair has a nonempty fibre, or that
all such fibres are empty. Existence or emptiness of the complete
sixteen-coordinate spatial projection remains unresolved by this result.
